\documentclass[12pt]{article}
\usepackage[T1]{fontenc}
\usepackage[utf8]{inputenc}
\usepackage{lmodern}
\usepackage{textcomp}
\usepackage{enumitem}
\usepackage{geometry}
\usepackage{graphicx}
\usepackage{amsmath, amssymb}
\geometry{margin=1in}
\begin{document}
\begin{center}
Sociology - Graduate Level Assessment 10 \
Total Marks: 40 \
Answer all questions.
\end{center}

\section*{Multiple Choice (10 marks)}
\begin{enumerate}[label=\arabic*.]
\item Which of the following is not identified as a new form of community?
\begin{enumerate}[label=(\alph*)]
    \item ethnic communities, based on shared identity and experiences of discrimination
    \item gay villages, which are formed in certain parts of large cities
    \item sociological communities, formed by unpopular lecturers
    \item virtual communities that exist only in cyberspace
\end{enumerate}
\item Scientific management involved:
\begin{enumerate}[label=(\alph*)]
    \item the subdivision of labour into small tasks
    \item the measurement and specification of work tasks
    \item motivation and rewards for productivity
    \item all of the above
\end{enumerate}
\item A survey should avoid asking:
\begin{enumerate}[label=(\alph*)]
    \item fixed-choice questions
    \item short questions
    \item leading questions
    \item funnelled questions
\end{enumerate}
\item Which of the following is not identified by Fulcher \& Scott as a criterion of community?
\begin{enumerate}[label=(\alph*)]
    \item a shared sense of identity and belonging together
    \item common activities involving all-round relationships
    \item a fixed geographical location
    \item collective action based on common interests
\end{enumerate}
\item 'Snowballing' is an example of:
\begin{enumerate}[label=(\alph*)]
    \item probability sampling
    \item non-probability sampling
    \item cluster sampling
    \item using the Christmas vacation constructively
\end{enumerate}
\end{enumerate}

\section*{True / False (10 marks)}
\textit{Answer True or False}
\begin{enumerate}[label=\arabic*.]
\setcounter{enumi}{5}
\item True or False: Goldthorpe defined the 'service class' as individuals in non-manual occupations who exercise authority on behalf of the state.
\item True or False: Traditional working class identity was primarily defined by collective aspirations to move into the middle class.
\item True or False: After the abolition of slavery, policymakers believed that former slaves would assimilate into urban society and its cultural norms.
\item True or False: Decarceration encompasses community alternatives to imprisonment and institutional care as a strategy for reducing incarceration rates.
\item True or False: Modernization theory posits that societies evolve from traditional structures to modern, industrialized forms of organization.
\end{enumerate}

\section*{Long Form Response (20 marks)}
\textit{Answer each question with a clear and complete explanation.}
\begin{enumerate}[label=\arabic*.]
\setcounter{enumi}{10}
\item Discuss the implications of globalization on the power dynamics of nation-states, specifically analyzing features that challenge or diminish their authority.
\item Discuss Blauner's (1964) analysis of alienation in the workplace, with a focus on car manufacturing in assembly plants. What factors contribute to employee alienation in this industry?
\end{enumerate}

\vfill
\noindent\textit{End of Paper}
\end{document}